\documentclass[10pt,a4paper]{amsart}
\usepackage[margin=19mm]{geometry}
\usepackage{amsmath,amssymb,mathtools}
\usepackage[scr=rsfs,cal=euler]{mathalfa}
\usepackage{xfrac}
\usepackage[unicode,hidelinks,bookmarks=false]{hyperref}
\newcommand{\ie}{i.e.\ }
\newcommand{\cf}{cf.\ }
\newcommand{\resp}{respectively\ }
\newcommand{\Subsection}{\subsection}
\providecommand{\nobreakdash}{\nobreak\hbox{-}}
\setlength{\emergencystretch}{2em}
\multlinegap0pt
\usepackage{xcolor}
\usepackage{algorithm}\usepackage{algpseudocode}
\usepackage{amssymb}
\usepackage[shortlabels]{enumitem}
\usepackage[normalem]{ulem}
\newtheorem{property}{Property}
\newtheorem{lemma}{Lemma}
\newtheorem{corollary}{Corollary}
\newtheorem{exam}{Exam}
\newcommand{\bN}{\mathbf{N}}
\newcommand{\bp}{\mathbf{p}}
\newcommand{\bq}{\mathbf{q}}
\newcommand{\defo}[1] {\emph{\textcolor{blue}{#1}}}
\renewcommand{\tt} {\mathtt}
\title{Decoding universal cycles for the $k$-subsets of an $n$-set}
\author{Daniel Gabri\'{c}, Wazed Imam, Lukas Janik-Jones, Joe Sawada, ~}
\date{Source: arXiv:2603.11934v1 (CC BY 4.0)}
\begin{document}
\maketitle

\noindent\emph{Source and licence.} Decoding universal cycles for the $k$-subsets of an $n$-set. Version arXiv:2603.11934v1 (CC BY 4.0), distributed by arXiv under the Creative Commons Attribution 4.0 International licence, http://creativecommons.org/licenses/by/4.0/, as recorded in the arXiv licence metadata for this exact version and shown on the abstract page https://arxiv.org/abs/2603.11934v1. Full licence notice: \texttt{licenses/arxiv-2603.11934-CC-BY-4.0.txt}.\par\smallskip

\noindent\textit{Editorial scope.} Counted unit: the article's lemma lem:lexorder (source0.tex lines 2194--2204). The article's complete original proof of it is delivered with it (source0.tex lines 2235--2282, one \texttt{proof} environment of 48 lines, which the article places away from the statement). No mathematics is written, repaired or re-derived: the statement and the proof are the article's own lines, in the article's own notation, and no retained line is substituted or re-flowed.  Retained uncounted as the source-local context the proof needs: lines 561--571; lines 2161--2165; lines 2208--2210.  Nothing was omitted from any retained span, so the package carries no omission record.  No retained line contains a citation, so the wrapper carries no bibliography and no bibliography entry.  No retained line contains a figure environment, an image-inclusion command or drawing code, so no image is delivered and the package declares no asset dependency.  Editorial formatting only, in addition: an AMS \texttt{amsart} preamble that restates the article's own declarations and macros used by the retained lines, namely the environments \texttt{property}, \texttt{lemma}, \texttt{corollary}, \texttt{exam} on separate counters, together with the standard packages those lines need.  The retained environments print in this document as Property 1, Lemma 1, Corollary 1, Exam 1 in order of appearance, whereas the article prints the same items with the numbering of its own full document.  The complete native source of the article as distributed in the arXiv e-print for this exact version is bundled unchanged as \texttt{sources/196244/source0.tex}.
%
\begin{property} \label{prop:first}
Let $n,k > 1$ and $k < w < kn$.  The first necklace in $\bN_k(n,w^\uparrow)$ is $\tt{1}^{n-j-1}\tt{x}\tt{k}^{j}$,
where $j=\lfloor \frac{w-n}{k-1}\rfloor$ and $\tt{x} = w-(n{-}j{-}1)-kj$; it has weight $w$ and is aperiodic.
\end{property}
%
\begin{property} \label{prop:last}
 Let $n,k > 1$ and $w < kn$.  The last two necklaces in  $\bN_k(n,w^\uparrow)$ are
 $(\tt{k{-}1})\tt{k}^{n-1}$ and $\tt{k}^n$.
\end{property}
%

\begin{lemma} \label{lem:max}
    Let $\beta = \tt{b}_1 \tt{b}_2 \cdots \tt{b}_{i-1} \cdot \tt{x} \cdot \tt{k}^{n-i}$ be a necklace in $\bN_k(n)$ and let $\tt{x}<k$. There exists a necklace $\beta_2 = \tt{b}_1 \tt{b}_2 \cdots \tt{b}_{i-1} \cdot \tt{(x+1)} \cdot \tt{k}^{n-i}$. In other words, increasing the last non-$k$ character of a necklace always produces a valid necklace.

    % Let $\beta$ be a necklace in $\bN_k(n)$ and let $\beta = \bp \bq$, where $\bq$ has length $i$. If the suffix $\bq$ is replaced with a maximal suffix, the resultant string $\bp \cdot \tt{k}^i$ is also a valid necklace in $\bN_k(n)$. If $\beta \neq \bp \cdot \tt{k}^i$, then $\bp \cdot \tt{k}^i$ occurs after $\beta$ lexicographically.
\end{lemma}

\begin{corollary} \label{cor:max2}
    For two necklaces $\beta_1$ and $\beta_2$ to be consecutive in $\bN_k(n)$, the suffix after the increasing index of $\beta_1$ must be a maximal suffix.
\end{corollary}

\begin{lemma} \label{lem:lexorder}
    Let $\beta_1$, $\beta_2$ be consecutive necklaces in the sequence of necklaces $\bN_k(n)$. There exists an index $1 \leq i\leq n$ such that
    \[
    \beta_1 = \tt{a}_1 \tt{a}_2\cdots \tt{a}_{i-1} \cdot \tt{x} \cdot \tt{a}_{i+1} \tt{a}_{i+2} \cdots \tt{a}_n
    \]\[
    \beta_2 = \tt{a}_1 \tt{a}_2\cdots \tt{a}_{i-1} \cdot \tt{y} \cdot \tt{b}_{i+1} \tt{b}_{i+2} \cdots \tt{b}_n
    \]
    where $\tt{x},\tt{y}\in \Sigma_k$ with $\tt{x}<\tt{y}$. $\beta_1$ and $\beta_2$ share the common prefix $\tt{a}_1 \tt{a}_2\cdots \tt{a}_{i-1}$, and the increasing index between them is $i$.
    
    % We refer to $\tt{a}_1 \tt{a}_2\cdots \tt{a}_{i-1}$ as the \defo{common prefix} of $\beta_1$ and $\beta_2$, and $i$ as the \defo{increasing index} of $\beta_1$ and $\beta_2$.
\end{lemma}

\begin{proof} 
%
We can divide this into a proof by cases by considering the location of the increasing index between $\beta_1$ and $\beta_2$. There are two cases that must be considered when searching for $\alpha$. \\

\noindent
\textbf{Case 1.} Suppose that the increasing index between $\beta_1$ and $\beta_2$ is within the suffix $\bp$. Let us consider the positions of $\bp$ and $\bq$ in $\bN_k(n,w^\uparrow)$:

\begin{enumerate}
    \item $\bp$: By Lemma \ref{lem:lexorder}, $\beta_1$ and $\beta_2$ must share a common prefix up until the increasing index, both necklaces must share the common prefix of $\bq$. This means that $\beta_1 = \bq \bp$, with weight at least $w$ by definition of $\alpha$. Furthermore, we know that $\bp$ is not a maximal suffix, as it contains an increasing index. So, $\beta_1$ contains the suffix $\bp$ and must exist in the sequence $\bN_k(n,w^\uparrow)$.

    \item $\bq$: Let $g$ be the length of $\bq$ and let $\beta_3$ be a necklace of the form $\bq \cdot \tt{k}^{n-g}$. By Lemma \ref{lem:max}, $\beta_3$ must come after $\beta_1$ in lexicographic order.

    Clearly, $\beta_3$ must be heavier than $\beta_1$, thus $\beta_3$ is also a necklace in $\bN_k(n,w^\uparrow)$. Lexicographically, as $\beta_1$ and $\beta_3$ have the same prefix, $\bq$, no necklace between the two of them in $\bN_k(n)$ can have a differing prefix. If any necklace from $\bN_k(n)$ between $\beta_1$ and $\beta_3$ is included in $\bN_k(n,w^\uparrow)$, it also has the prefix $\bq$. Since $\beta_3$ is in $\bN_k(n,w^\uparrow)$, this means the the first necklace in $\bN_k(n,w^\uparrow)$ lexicographically larger than $\beta_1$ must have the prefix $\bq$.
\end{enumerate}

Therefore, in this case we can derive the location of $\alpha$ in $\bN_k(n,w^\uparrow)$ using $\bN_k(n)$. If $\alpha = \bp \bq$ in $\bN_k(n)$, where $\bp$ is the suffix of some necklace $\beta_1$ and $\bq$ is the prefix of some necklace $\beta_2$, then in $\bN_k(n,w^\uparrow)$ $\beta_1$ also appears in $\bN_k(n,w^\uparrow)$, and $\bq$ is found in the subsequent necklace. \\

\noindent
\textbf{Case 2.} Suppose that the increasing index between $\beta_1$ and $\beta_2$ is within the prefix $\bq$. Let us consider the positions of $\bp$ and $\bq$ in this case as well:

\begin{enumerate}
    \item $\bq$: By Lemma \ref{lem:lexorder}, $\beta_1$ must have a prefix lexicographically smaller than $\bq$. Additionally, Corollary \ref{cor:max2} tells us that $\bp$ must be a maximal suffix of the form $\tt{k}^g$. Similarly to the previous case, we can use Lemma \ref{lem:max} to construct a necklace $\beta_3$ which contains the prefix $\bq$ and the maximal suffix $\bp$.
    
    Let $\beta_3 = \bq \bp$. As $\alpha=\bp \bq$ has a weight of at least $w$, $\beta_3$ must also have a weight of at least $w$.Once again, we can observe that any necklace between $\beta_2$ and $\beta_3$ in $\bN_k(n)$ must contain the prefix $\bq$. Therefore, as $\beta_3$ is in $\bN_k(n,w^\uparrow)$, the necklace following $\beta_1$ in $\bN_k(n,w^\uparrow)$ has the prefix $\bq$.

    \item $\bp$: Unlike the first case, it is not certain that $\beta_1$ has a weight of $w$ or greater and may not be present in $\bN_k(n,w^\uparrow)$. However, suppose that there is some necklace in $\bN_k(n,w^\uparrow)$ smaller than $\beta_1$ lexicographically, but does not contain the suffix $\bp$. Since $\bp$ is a maximal suffix, we can construct another necklace, $\beta_0$, which has the same prefix as this necklace, but contains the maximal suffix $\bp$. 

    Let $\beta_0$ be a necklace that has a prefix lexicographically smaller than $\bq$, and containing the maximal suffix $\bp$. By Lemma \ref{lem:max}, this necklace exists. Note that $\beta_0$ is both heavier than the supposed previous necklace, and comes between that necklace and $\beta_2$ lexicographically, as $\beta_2$ necessarily contains a larger prefix. In other words, if any necklace exists before $\beta_1$ in $\bN_k(n,w^\uparrow)$ that does not contain the suffix $\bp$, there is always another necklace $\beta_0$ that contains $\bp$ as a suffix that comes before $\beta_2$. 

    \begin{exam}
    Suppose $n=4$, $k=3$, and $w=7$. Let $\alpha = \tt{\textcolor{blue}{33}\textcolor{red}{11}}$ $\beta_1 = \tt{0322}$, and $\beta_2 = \tt{\textcolor{red}{11}23}$. $\beta_1$ comes before $\beta_2$ in $\bN_k(n,w^\uparrow)$. As $\tt{\textcolor{blue}{33}}$ is a max suffix, we can construct the following necklace with the same prefix as $\beta_1$:
    %
    \[ \beta_0 = \tt{03\textcolor{blue}{33}} \]
    %
    This necklace appears between $\beta_1$ and $\beta_2$, and contains the maximal suffix found in $\alpha$. This is true regardless of the prefix of $\beta_1$. In this case, the sequence becomes:
    
    \[ \tt{\cdots, 0322, \ldots, 03\textcolor{blue}{33}, \ldots, \textcolor{red}{11}23}, \ldots \]
    \end{exam}
    
    If no necklace exists in $\bN_k(n,w^\uparrow)$ before $\beta_2$, then recall Property \ref{prop:last}. When considered cyclically the previous necklaces in $\bN_k(n,w^\uparrow)$ are $(k-1) \tt{k}^{n-1}$ and $k^n$. Since the length of $\bp$, $g$, is less than $n$, the previous suffix would once again be $\bp = \tt{k}^g$. In all cases, the suffix preceding $\beta_2$ is $\bp$.

\end{enumerate}

% Step three: put it together

Therefore, in both cases we derive the position of $\alpha = \bp \bq$ in $\bN_k(n,w^\uparrow)$ using the position of $\alpha$ in $\bN_k(n)$. Regardless of the increasing index of $\alpha$, the prefix $\bq$ is always contained in the first necklace in $\bN_k(n)$ after $\beta_1$ that has a weight of at least $w$.
%
\end{proof}

\end{document}
